% Checked addition; independently audited by root on 2026-11-20.
\subsubsection*{Statement B.1f (Proved): finite-group detection over an arbitrary field}
Let $k$ be a field of characteristic $p>0$, let $P$ be a commutative
group scheme of finite type over $k$, and let $Q\subset P$ be a closed
subgroup scheme.  If every finite $k$-subgroup scheme of $P$ is contained
in $Q$, then $Q=P$.
\begin{proof}
For each $r\geq 1$, the relative Frobenius $F^r:P\to P^{(p^r)}$ is
finite.  On an affine chart this follows because a finitely generated
$k$-algebra is a finite module over the subalgebra generated by $k$
and all its $p^r$-th powers.  Thus $\ker F^r$ is a finite $k$-group
scheme and is contained in $Q$.

After extending to $\bar k$, put
$R=\mathcal O_{P_{\bar k},e}$ and let $\mathfrak m$ be its maximal
ideal.  The ideal of $\ker F^r$ at $e$ is generated by $p^r$-th powers
of elements of $\mathfrak m$.  The ideal $I$ of $Q_{\bar k}$ therefore
satisfies
\[
 I\subset \mathfrak m^{p^r}\quad\text{for every }r,
 \qquad I\subset\bigcap_{n\geq1}\mathfrak m^n=0.
\]
Here the last equality is the Krull intersection theorem.  Consequently
$Q_{\bar k}$ contains an open neighbourhood of the identity, and is
an open and closed subgroup of $P_{\bar k}$.  In particular it contains
the entire scheme-theoretic identity component.

We show that it meets every geometric component.  Set
$G=(P_{\bar k})_{\mathrm{red}}^0$.  The structure theorem for smooth
commutative algebraic groups gives a smooth connected unipotent subgroup
$U\subset G$ with $G/U$ semiabelian
\cite[Theorem 8.27 and Theorem 16.13]{MilneAG2022B}.
There is an integer $a$ such that $[p^a]U=0$: embed $U$ into a group
of upper triangular unipotent matrices and take $p^a$ at least the
matrix size.  For any $x\in P(\bar k)$, choose $m>0$ such that
$mx\in G(\bar k)$.  Multiplication by $m$ is surjective on the
semiabelian variety $G/U$.  We can therefore choose $y\in G(\bar k)$
with $m(x-y)\in U(\bar k)$.  The point $t=x-y$ lies in the component
of $x$ and is killed by $mp^a$.

The coordinates of $t$ are algebraic over $k$.  For sufficiently large
$r$, the point $t'=F^r(t)$ of $P^{(p^r)}$ is defined over $k^{\mathrm{sep}}$.
The finite Galois orbit of $t'$ generates a finite group, since the
group is commutative and every generator has the same finite-order
bound.  This Galois-stable set of geometric points defines a finite
\'etale $k$-subgroup $D\subset P^{(p^r)}$.  The inverse image
$H=(F^r)^{-1}(D)$ is a finite $k$-subgroup of $P$ and contains $t$
after extension to $\bar k$.  By hypothesis $H\subset Q$, so
$Q_{\bar k}$ meets the component of $x$.  Being open and closed and
containing the identity component, it therefore equals $P_{\bar k}$.
Faithfully flat descent gives $Q=P$.
\end{proof}

\subsubsection*{Statement B.1g (Proved): normal source and regular target over any field}
Let $k$ be a field of characteristic $p>0$, and let $f:X\to Y$ be a
proper birational morphism of proper geometrically integral
$k$-schemes.  If $X$ is normal and $Y$ is regular, then
\[
 f^*:\operatorname{Pic}^{\tau}_{Y/k}
       \xrightarrow{\ \sim\ }\operatorname{Pic}^{\tau}_{X/k}.
\]
Geometric normality of $X$ or $Y$ is not required in this statement.
\begin{proof}
The two Picard torsion components are group schemes of finite type
\cite[Corollary 4.18.3, Proposition 6.12 and Remarks 6.14 and 6.19]{Kleiman2005}.
Normality of $Y$ gives $f_*\mathcal O_X=\mathcal O_Y$ universally over
$k$.  The projection formula consequently makes $f^*$ a monomorphism;
a monomorphism of algebraic groups is a closed immersion
\cite[Theorem 5.34]{MilneAG2022B}.  Write its image as $Q\subset P$.

Geometric integrality implies that $X$ has a nonempty smooth locus.
A point in this locus is defined over a finite separable extension
of $k$.  We may make this extension, since normality and regularity
are preserved and an isomorphism of the Picard schemes descends.
We then have rational points of $X$ and $Y$ compatible with $f$.
For every finite commutative $k$-group $H$, Raynaud's correspondence
and the sections give
\[
 \operatorname{Hom}(H,\operatorname{Pic}_{Z/k})
   =H^1(Z,H^D)/H^1(k,H^D),\qquad Z=X,Y,
\]
as in B.1e.  Finite-flat torsor purity B.1b makes these groups
isomorphic.  Each homomorphism from $H$ lands in $\operatorname{Pic}^
{\tau}$, because its component-group image is torsion.  In particular
every finite $k$-subgroup scheme of $P$ factors through $Q$.
Statement B.1f now gives $Q=P$.
\end{proof}
